\chapter{Introduction}
\label{ch:intro}

Start writing your notes here. This chapter is a placeholder that
demonstrates the environments provided by the class, so you can see how
they render before replacing the content with your own.

\section{Definitions and Results}

The class (via \texttt{amsthm}) gives you numbered environments for the
usual building blocks of a maths course:

\begin{definition}
\label{def:example}
A \textbf{definition} environment looks like this. Numbers follow the
chapter, e.g. Definition~\ref{def:example}.
\end{definition}

\begin{proposition}
\label{prop:example}
A \textbf{proposition} environment looks like this.
\end{proposition}

\begin{proof}
And here is how a proof looks, using \texttt{amsthm}'s \texttt{proof}
environment (inherited via \texttt{amsthm}, which \texttt{Base.cls}
loads).
\end{proof}

\begin{theorem}
\label{thm:example}
A \textbf{theorem} environment.
\end{theorem}

\begin{lemma}
\label{lem:example}
A \textbf{lemma} environment.
\end{lemma}

\begin{corollary}
\label{cor:example}
A \textbf{corollary} environment, sharing the same counter as theorems.
\end{corollary}

\begin{example}
\label{ex:example}
An \textbf{example} environment, useful for worked problems.
\end{example}

\begin{remark}
A \textbf{remark} environment (numbered per section instead of per
chapter).
\end{remark}

\section{Citing Sources}

You can cite references from \texttt{refs.bib}, e.g.\ \cite{knuth1986texbook}.

% Note: the Exercises environment below inserts its own numbered
% "Exercises" \section heading automatically -- don't add one by hand.
\begin{Exercises}
\item This is the first exercise in the chapter.
\item \Hint This is a hint for the second exercise.
\end{Exercises}
